\section*{Chapitre \ref{ch:nw:buying-connectivity} --- Achat de connectivité}

\begin{solution}{exo:nw:buying-connectivity:1}
MTBF $8\,766/2 = \qty{4383}{\hour}$, MTTR \qty{4}{\hour}~: disponibilité $4\,383/4\,387 = 99.9088~\%$, indisponibilité $8\,766 \times 60 \times 0.000912 = 480$ minutes par
an (huit heures).
\end{solution}

\begin{solution}{exo:nw:buying-connectivity:2}
$9\,000/10 = 900$~USD par mégabit et par mois à 10~mégabits~; $44\,000/200 = 220$~USD à 200~mégabits.
\end{solution}

\begin{solution}{exo:nw:buying-connectivity:3}
En dessous de 95~\% de $30 \times 1\,440 = 43\,200$ minutes~: plus de 2\,160 minutes, soit 36 heures, au cours du mois.
\end{solution}

\begin{solution}{exo:nw:buying-connectivity:4}
La paire est indisponible $u^2 + c$~: avec $u = 0.00091$, le terme indépendant vaut $8.3 \times 10^{-7}$, négligeable, et celui de la gaine, $c = 6/8\,766 = 0.00068$
est du même ordre que le $u$ d'un circuit. Le second circuit supprime les pannes indépendantes, qui n'ont jamais été le problème, et conserve la
cause commune.
\end{solution}

\begin{solution}{exo:nw:buying-connectivity:5}
L'entrée de l'immeuble ou la rue voisine~; la \omterm{def:nw:colocation-products-and-how-they-are-sold:mmr}{salle d'interconnexion} (meet-me room) ou la cage de l'opérateur~; un panneau de brassage ou un commutateur à chaque extrémité~; l'alimentation électrique des
baies~; le fournisseur lui-même (faute logicielle, erreur de maintenance, faillite) si les deux viennent du même~; une fibre de gros acheminée sous deux marques commerciales.
\end{solution}

\begin{solution}{exo:nw:buying-connectivity:6}
Lorsque l'entreprise a besoin de nombreuses plateformes et sources de données à une latence modérée, ou d'un accès rapide à une nouvelle plateforme, et ne souhaite pas construire ni exploiter
des circuits vers chacune. Non lorsque l'edge de la stratégie est la latence vers cette plateforme~: le réseau propre de l'extranet et sa capacité partagée
se trouvent sur le trajet, et l'entreprise ne peut pas les contrôler.
\end{solution}

\begin{solution}{exo:nw:buying-connectivity:7}
Résoudre $u^2 + c - u^2 c = u$ avec les circuits du modèle donne $c = 0.000911$, soit 0.67 événement de gaine par an de douze heures~: à ce rythme, le second
circuit dans la même gaine n'apporte strictement rien.
\end{solution}

\begin{solution}{exo:nw:buying-connectivity:8}
Le crédit est une fraction de la redevance du mois pour le service défaillant (10~\% au premier seuil), imputée sur les factures suivantes, et c'est le
seul recours~; la perte de trading de l'entreprise pendant la panne n'est pas couverte. Et 99.99~\% d'un mois laisse encore quatre minutes de panne par mois sans
aucun crédit du tout.
\end{solution}

\begin{solution}{pb:nw:buying-connectivity:1}
\textbf{Partie I.}
\begin{enumerate}
\item 99.9088~\%, 479.6 minutes par an.
\item $12 \times (11\,000 + 600) = 139\,200$~USD par an.
\item $479.6 \times 2\,000 \approx 959\,000$~USD.
\item Rien~: quatre heures dans un mois laissent une disponibilité de 99.44~\%, au-dessus du seuil de 95~\%.
\end{enumerate}
\textbf{Partie II.}
\begin{enumerate}[resume]
\item 99.9315~\%, 360.4 minutes par an.
\item 360.0 des 360.4 minutes, soit un demi-événement par an de douze heures.
\item 278\,400~USD par an~; elle évite environ 238\,000~USD de perte espérée par rapport à un circuit, soit moins que son surcoût de 139\,200, plus le risque qu'elle conserve.
\item 1\,992 minutes, soit 33 heures~: pire que les 1\,541 d'un circuit, car un événement de gaine dure douze heures.
\end{enumerate}
\textbf{Partie III.}
\begin{enumerate}[resume]
\item 99.99992~\%, 0.44 minutes par an.
\item 41\,760~USD par an de plus que la gaine~; elle économise 720\,000~USD de perte espérée.
\item Un circuit coûte environ 5.50~USD, la paire en gaine environ 2\,560~USD, la paire diversifiée rien, selon les barèmes de leurs niveaux de service.
\item Un élément commun oublié dans l'étude (entrée de l'immeuble, alimentation, systèmes propres du fournisseur), une catastrophe régionale, ou une panne à l'extrémité de la
  plateforme, qu'aucun circuit ne corrige (chapitre~28).
\end{enumerate}
\textbf{Partie IV.}
\begin{enumerate}[resume]
\item \emph{Résultat nommé}~: deux circuits dans une même gaine sont disponibles 99.93~\% du temps (360 minutes d'indisponibilité par an, 721\,000~USD de perte espérée)~; deux
  circuits diversifiés à 99.99992~\% (0.4 minutes)~; les \omterm{def:nw:buying-connectivity:sla}{crédits de service} gagnés par la paire en gaine représentent environ 2\,560~USD par an, soit moins de 0.4~\% de sa perte espérée.
\item La fréquence et la durée des événements de mode commun~: toute la différence entre les conceptions s'y joue.
\item Demander les plans de tracé des deux circuits jusqu'à la rue et l'entrée de l'immeuble, les confronter entre eux et aux
  infrastructures partagées des fournisseurs, et tester en coupant un circuit.
\item La latence et sa stabilité (pas seulement la disponibilité), le temps de réparation, le préavis des maintenances et les engagements de diversité du fournisseur.
\item Deux circuits actifs valent mieux~: le trafic est réparti ou dupliqué, la détection de panne est immédiate, et le secours est connu pour fonctionner~; une configuration actif-secours
  exige des tests de bascule réguliers pour inspirer confiance.
\item Dans sa section sur la résilience~: la conception de chaque connexion, sa disponibilité et son coût, et la perte qu'elle protège.
\item Une protection contre les pannes indépendantes, qui ne représentaient qu'une faible part du risque~; elle ne protège pas contre la gaine.
\item Deux éléments qui tombent en panne séparément.
\end{enumerate}
\end{solution}

\begin{solution}{iq:nw:buying-connectivity:1}
La clause de performance du contrat (disponibilité, parfois latence et temps de réparation), la façon dont elle est mesurée, et ce que le fournisseur doit
lorsqu'il manque à ses engagements. Un crédit indemnise la redevance du service non fourni, généralement une fraction de la redevance mensuelle~; il n'indemnise pas les
pertes du client.
\iqlookfor{Période et méthode de mesure~; crédits exprimés en fraction des redevances~; recours unique~; exclusions.}
\end{solution}

\begin{solution}{iq:nw:buying-connectivity:2}
Par tout ce qu'ils partagent~: une gaine, une entrée d'immeuble, une salle, un panneau de brassage, l'alimentation, un opérateur de gros commun aux deux, ou le côté de la plateforme dans la connexion~;
et par des causes communes~: une catastrophe régionale, une configuration déployée sur les deux.
\iqlookfor{Modes communs physiques et logiques~; vérification par les plans de tracé et par des tests.}
\end{solution}

\begin{solution}{iq:nw:buying-connectivity:3}
Pour la largeur (breadth) et la rapidité de mise en service (de nombreuses plateformes et flux par une seule connexion, une nouvelle plateforme en quelques semaines), lorsque la stratégie n'est pas critique en latence sur
cette plateforme~; pas pour les plateformes où l'entreprise se bat de vitesse.
\iqlookfor{Latence contre commodité~; le réseau partagé du fournisseur~; accords avec les plateformes et frais.}
\end{solution}

\begin{solution}{iq:nw:buying-connectivity:4}
La disponibilité est $\mathrm{MTBF}/(\mathrm{MTBF}+\mathrm{MTTR})$~; avec une indisponibilité $u$ chacun, deux circuits indépendants sont simultanément en panne pendant $u^2$ du temps,
par exemple $0.001^2 = 10^{-6}$~; mais tout mode commun $c$ s'ajoute directement et domine en général.
\iqlookfor{La formule~; les carrés pour l'indépendance~; le mode commun comme risque réel.}
\end{solution}

\begin{solution}{iq:nw:buying-connectivity:5}
Pour le temps, et non la bande passante~: quelques mégabits qui arrivent des centaines de microsecondes ou de millisecondes plus tôt portent les messages qui tranchent les courses de vitesse, et
leur valeur dépasse le prix~; tout le reste passe par la fibre.
\iqlookfor{Valeur de la latence~; petits messages~; secours par fibre.}
\end{solution}

\begin{solution}{iq:nw:buying-connectivity:6}
Deux chemins diversifiés depuis le bureau vers chaque ville (fournisseurs, rues, entrées et alimentations différents), des connexions doubles à chaque plateforme depuis deux
points diversifiés, les deux actifs~; vérifier les plans de tracé, tester régulièrement la bascule, surveiller la latence et les pertes sur chaque chemin, et calculer la disponibilité de la conception avec
des hypothèses explicites sur les modes communs~; cinq neuf représentent environ cinq minutes par an, ce que seule la diversité permet d'obtenir.
\iqlookfor{Diversité de bout en bout~; actif-actif~; tests~; l'arithmétique de 99.999~\%.}
\end{solution}
